\documentclass[11pt]{article}

\usepackage[T1]{fontenc}
\usepackage[utf8]{inputenc}
\usepackage{lmodern}
\usepackage[margin=1.04in]{geometry}
\usepackage{amsmath,amssymb,mathtools,amsthm}
\usepackage{aliascnt}
\usepackage{mathrsfs}
\usepackage{microtype}
\usepackage{enumitem}
\usepackage{xcolor}
\usepackage{booktabs}
\usepackage{url}
\usepackage{hyperref}
\usepackage[capitalize,noabbrev]{cleveref}

\hypersetup{
  colorlinks=true,
  linkcolor=blue!45!black,
  citecolor=green!35!black,
  urlcolor=blue!55!black,
  pdftitle={Simplex Exponential Periods: Reductions and Open Comparison Problems},
  pdfauthor={Christopher D. Long; Antoine-Auguste Le Blanc},
  pdfsubject={Candidate simplex nonvanishing program and audited reductions},
  pdfkeywords={exponential periods, candidate draft, comparison, open problems}
}

\numberwithin{equation}{section}
\setlength{\emergencystretch}{3em}

\newtheorem{theorem}{Theorem}[section]

\newaliascnt{proposition}{theorem}
\newtheorem{proposition}[proposition]{Proposition}
\aliascntresetthe{proposition}

\newaliascnt{lemma}{theorem}
\newtheorem{lemma}[lemma]{Lemma}
\aliascntresetthe{lemma}

\newaliascnt{corollary}{theorem}
\newtheorem{corollary}[corollary]{Corollary}
\aliascntresetthe{corollary}

\theoremstyle{definition}
\newaliascnt{definition}{theorem}
\newtheorem{definition}[definition]{Definition}
\aliascntresetthe{definition}

\newaliascnt{example}{theorem}
\newtheorem{example}[example]{Example}
\aliascntresetthe{example}

\theoremstyle{remark}
\newaliascnt{remark}{theorem}
\newtheorem{remark}[remark]{Remark}
\aliascntresetthe{remark}

\newcommand{\C}{\mathbb C}
\newcommand{\Q}{\mathbb Q}
\newcommand{\Qbar}{\overline{\mathbb Q}}
\newcommand{\R}{\mathbb R}
\newcommand{\Z}{\mathbb Z}
\newcommand{\N}{\mathbb N}
\newcommand{\Aone}{\mathbb A^1}
\newcommand{\Atwo}{\mathbb A^2}
\newcommand{\Gm}{\mathbb G_m}
\newcommand{\Spec}{\operatorname{Spec}}
\newcommand{\Span}{\operatorname{span}}
\newcommand{\trdeg}{\operatorname{trdeg}}
\newcommand{\per}{\operatorname{per}}
\newcommand{\ord}{\operatorname{ord}}
\newcommand{\End}{\operatorname{End}}
\newcommand{\Hom}{\operatorname{Hom}}
\newcommand{\Ext}{\operatorname{Ext}}
\newcommand{\Gal}{\operatorname{Gal}}
\newcommand{\Lie}{\operatorname{Lie}}
\newcommand{\Frac}{\operatorname{Frac}}
\newcommand{\Sym}{\operatorname{Sym}}
\newcommand{\SL}{\operatorname{SL}}
\newcommand{\GL}{\operatorname{GL}}
\newcommand{\Pic}{\operatorname{Pic}}
\newcommand{\Div}{\operatorname{Div}}
\newcommand{\Res}{\operatorname{Res}}
\newcommand{\Tr}{\operatorname{Tr}}
\newcommand{\Crit}{\operatorname{Crit}}
\newcommand{\supp}{\operatorname{supp}}
\newcommand{\rank}{\operatorname{rank}}
\newcommand{\Ecal}{\mathcal E}
\newcommand{\Fcal}{\mathcal F}
\newcommand{\Hcal}{\mathcal H}
\newcommand{\Lcal}{\mathcal L}
\newcommand{\Mcal}{\mathcal M}
\newcommand{\Ocal}{\mathcal O}
\newcommand{\Pcal}{\mathcal P}
\newcommand{\Rcal}{\mathcal R}
\newcommand{\Vcal}{\mathcal V}
\newcommand{\dd}{\,d}
\newcommand{\KZ}{\mathrm{KZ}}

\title{Simplex Exponential Periods: Reductions and Open Comparison Problems}
\author{Christopher D. Long \and Antoine-Auguste Le Blanc\thanks{Le Blanc cryptographic author identifier: \texttt{b6048125\allowbreak{}30b3535b\allowbreak{}23d6dbb0\allowbreak{}f5abe32d\allowbreak{}1ed1a9de\allowbreak{}35fb4f7d\allowbreak{}04d62cff\allowbreak{}08ea19c8}}}
\date{October 2026}


\newcommand{\Hq}{H_q^2}
\newcommand{\M}{\mathscr M}
\newcommand{\Ncal}{\mathcal N}
\newcommand{\Tcal}{\mathcal T}
\newcommand{\A}{\mathbb A}

\newcommand{\PGL}{\operatorname{PGL}}
\newcommand{\wt}{\operatorname{wt}}
\newcommand{\Bcal}{\mathcal B}

\newtheorem{conjecture}[theorem]{Conjecture}
\begin{document}
\maketitle

\begin{center}
\small
\texttt{galizur@gmail.com}
\end{center}

\begin{abstract}
This candidate draft separates the unrestricted simplex exponential
nonvanishing problem from the controlled higher-dimensional families proved
in Paper VI. We record the E-function reduction, a conjugate-real case,
fixed-phase effective exceptional sets, and an exact positive-amplitude
counterexample to a tempting induction. We formulate the actual-chain
visibility and higher-asymptotic-coefficient problems that remain.
No unrestricted simplex nonvanishing theorem, general numerical period
comparison, or new dimension-wide Factorial Conjecture is proved here.
\end{abstract}

\begin{center}
\textbf{Candidate research draft; not a numbered paper in the series.}
\end{center}
\tableofcontents

\section{The problem and its normalization}
Let
\[
 \Delta_d=\{t\in\R_{\ge0}^{d}:t_1+\cdots+t_d\le1\},\qquad
 F_Q(z)=\int_{\Delta_d}e^{zQ(t)}\,dt,
 \quad Q\in\Qbar[t_1,\ldots,t_d].
\]
The measure is coordinate Lebesgue measure; $F_Q(0)=1/d!$.
For $d=0$ the integral is the exponential of the constant phase.

\begin{conjecture}[Simplex exponential nonvanishing]\label{conj:simplex}
For every $d\ge1$ and every polynomial $Q\in\Qbar[t_1,\ldots,t_d]$,
$F_Q(1)\ne0$.
\end{conjecture}
Absorbing a nonzero algebraic scale into $Q$ makes this equivalent to
$F_Q(\xi)\ne0$ at every nonzero algebraic $\xi$. Positivity settles
real-valued $Q$. For complex coefficients positivity alone says nothing
about cancellation.

The one-dimensional constant-amplitude result is covered by Paper I
\cite{PaperI}. Paper V \cite{PaperV} treats explicit planar comparison
systems; Paper VI \cite{PaperVI} converts controlled weighted planar
periods into simplex integrals in every ambient dimension. Those statements
do not imply \cref{conj:simplex}: the phase obtained by planar slicing
still depends on only two active coordinates.

\section{E-functions and a conjugate-real case}
\begin{proposition}\label{prop:E}
For each fixed algebraic polynomial $Q$, $F_Q$ is an $E$-function.
Its Taylor coefficients and a nonzero linear differential annihilator over
$\Qbar(z)$ can be computed from $Q$.
\end{proposition}
\begin{proof}
Use the polynomial cube parametrization
\[
 t_j=u_j\prod_{i<j}(1-u_i)\quad(1\le j\le d),\qquad
 J(u)=\prod_{i=1}^{d-1}(1-u_i)^{d-i}.
\]
It is a diffeomorphism on the interiors, apart from measure-zero boundary
identifications. Thus
\[
 F_Q(z)=\int_{[0,1]^d}J(u)e^{z\widetilde Q(u)}\,du
       =\sum_{m\ge0}a_m\frac{z^m}{m!},\qquad
 a_m=\int_{[0,1]^d}J(u)\widetilde Q(u)^m\,du.
\]
All $a_m$ lie in one number field. Choose a positive integer $A$ clearing
the algebraic coefficient denominators of $\widetilde Q$. Monomial
integration divides by a product of $d$ positive integers bounded by
$Cm+C_0$. Therefore
\[
 A^m\operatorname{lcm}(1,\ldots,Cm+C_0)^d a_m
\]
is an algebraic integer. These clearing integers, and the common clearing
integer for $a_0,\ldots,a_m$, grow at most exponentially. Every conjugate
of $a_m$ is bounded by an exponential, by integration over the fixed cube.
These are the arithmetic $E$-function bounds.

The integrand is holonomic: $e^{z\widetilde Q}$ is annihilated by its
first-order polynomial connection, and multiplication by $J$ preserves
holonomicity. Closure under definite integration, with boundary terms
included, gives a univariate holonomic function. Iterated integration over
constant endpoints $0,1$ can be carried out by holonomic elimination;
see the holonomic integration framework in \cite{Zeilberger}. Endpoint
terms remain holonomic, so an inhomogeneous telescoper can be homogenized.
This gives a differential annihilator. Monomial integration computes
any required finite list of initial coefficients. No claim that a chosen
annihilator has its only finite singularity at zero is made.
\end{proof}

\begin{lemma}[Algebraic zeros survive coefficient conjugation]\label{lem:conjugation}
Let $F$ be an $E$-function with algebraic coefficients. If $F(1)=0$, then
$F^\tau(1)=0$ for every coefficient embedding $\tau$ into $\C$.
\end{lemma}
\begin{proof}
The proof of Beukers' Proposition~4.1 \cite{Beukers} gives this statement
explicitly: every local solution of the minimal differential equation
vanishes at the algebraic zero, and the same holds for its coefficient
conjugates. The conjugate $E$-function is one of those solutions. An
embedding of its coefficient field extends to $\Qbar$.
This is not the unjustified assertion that a field automorphism commutes
with an infinite analytic sum.
\end{proof}

\begin{corollary}[Conjugate-real phases]\label{cor:real}
Suppose the coefficients of $Q$ lie in a number field $K$ having an
embedding $\tau:K\hookrightarrow\C$ for which $Q^\tau$ has real
coefficients. Then $F_Q(1)\ne0$.
\end{corollary}
\begin{proof}
The simplex monomial moments are rational, so coefficient conjugation of
the Taylor series gives $(F_Q)^\tau=F_{Q^\tau}$.
If $F_Q(1)$ vanished, \cref{lem:conjugation} would contradict
$F_{Q^\tau}(1)>0$.
\end{proof}
This applies in particular when the whole coefficient field has a real
embedding, not merely when each coefficient separately has some real
conjugate. For a scaled integral, apply it to the coefficients of $\xi Q$
together.

\section{Fixed-phase exceptional sets are finite, not necessarily empty}
\begin{proposition}\label{prop:exceptional}
For fixed $Q$, the set of algebraic zeros of $F_Q$ is finite and can be
determined effectively from differential and initial-value data.
\end{proposition}
\begin{proof}
An algebraic entire $E$-function is a polynomial. If $F_Q$ is a polynomial,
its algebraic zeros are computable and finite since $F_Q(0)=1/d!\ne0$.
Otherwise apply Adamczewski--Rivoal's algorithm, Theorem~2 of \cite{AR},
to compute all algebraic arguments where $F_Q$ takes an algebraic value,
together with the values; retain those equal to zero. The required
annihilator and initial coefficients are supplied by \cref{prop:E}.
\end{proof}
The algorithm may be expensive. This proposition neither makes the
exceptional set empty nor proves that the specific argument $1$ is absent.
In particular an ordinary-point lifting statement cannot be applied at an
unexamined apparent singularity of an arbitrarily chosen scalar operator.

\section{Why positive-amplitude induction fails}
\begin{proposition}\label{prop:counterexample}
There is a strictly positive rational polynomial $p$ on $[0,1]$ with
$\int_0^1 p(t)e^{5it}\,dt=0$.
\end{proposition}
\begin{proof}
Set $g(t)=t^2(1-t)^2$ and $p=g''+25g$. Direct expansion gives
\[
 p(t)=25\left((t-\tfrac12)^2-\tfrac1{100}\right)^2+\frac{14}{25}>0.
\]
On the other hand,
\[
 \frac{d}{dt}\bigl((g'-5ig)e^{5it}\bigr)=p(t)e^{5it}.
\]
Both $g$ and $g'$ vanish at $0$ and $1$, so the integral is zero.
\end{proof}
This is an exact polynomial twisted-Stokes cancellation. It is not a
counterexample to \cref{conj:simplex}, which has constant amplitude and
the standard simplex measure. It does prohibit an induction that replaces
the simplex by an arbitrary positive polynomial weight and assumes
nonvanishing for that broader class. Integrating out a coordinate often
introduces an amplitude; that step must retain its specific algebraic
structure.

\section{Two distinct higher-dimensional extensions}
\subsection{Ambient-dimension extension}
For a polynomial phase in only $u=t_1,v=t_2$ and $n\ge3$,
\[
 \int_{\Delta_{n-1}}e^{Q(t_1,t_2)}\,dt
 =\frac1{(n-3)!}\int_{\Delta_2}(1-u-v)^{n-3}e^{Q(u,v)}\,du\,dv.
\]
Here $\Delta_{n-1}$ can equivalently be represented by $n$ barycentric
coordinates summing to one. Paper VI proves nonvanishing for two explicit
families of these weights and phases. One uses standard twisted remainder
basis elements and permits every nonzero algebraic scale; the other uses
a characteristic-five nonexactness theorem for all exponents.

\subsection{Genuine phase-dimension extension}
For an arbitrary $Q(t_1,\ldots,t_d)$ the interior twisted module can have
several constituents. Its relative extension by face systems and the
actual simplex class must both be controlled. A larger tensor category or
a formal comparison statement by itself does not ensure that this
particular scalar period survives. The substantive questions are:
\begin{enumerate}[leftmargin=2em]
\item Which irreducible interior constituent is detected by the actual
simplex integration map modulo the linear face span?
\item Can a scalar relation created at an algebraic parameter be lifted
in a closed nonsingular $E$-function system containing that constituent?
\item When several constituents occur, what prevents cancellation of their
contributions to the selected constant-amplitude coordinate?
\end{enumerate}
The elementary irreducible-module criterion in Paper V answers these
questions only after irreducibility and nonzero detection have been proved.
It is not a replacement for proving them.

\section{A potentially weaker route to factorial nonvanishing}
For the factorial functional of Paper VI, a flat highest form $S^D$ gives
\[
 \frac{\Lcal_n(f^m)}{\Gamma(Dm+n)}\longrightarrow F_Q(1),
 \qquad Q=f_{D-1}/D\bigm|_{S=1}.
\]
If this leading coefficient vanishes, no counterexample to the Factorial
Conjecture follows. Higher asymptotic coefficients involve polynomial
amplitudes multiplying $e^Q$, with those amplitudes determined by the
lower homogeneous layers and the gamma-ratio expansion.

\begin{conjecture}[Asymptotic visibility target]\label{conj:visibility}
For every nonzero polynomial with highest homogeneous form $S^D$, some
coefficient in a justified factorial-moment asymptotic expansion is
nonzero, or a separately controlled exponentially smaller contribution
prevents all positive moments from vanishing.
\end{conjecture}
This is a research target, not a theorem or a claimed formal equivalence.
An expansion with all power-series coefficients zero can still have
nonzero exponentially small terms. Any proposed proof must control its
remainder, not only manipulate a formal expansion.

\section{Audit queue and relation to the series}
The recovered FC(3) handoffs mention quadratic and cubic triangle
nonvanishing drafts. Their complete original proof sources were not part
of this integration. In particular, a manuscript claim about arbitrary
cubic triangle phases and degenerations should not be promoted on the
strength of a context summary. A dedicated audit must inspect the
relation-row exclusion, the relevant local lattices, and value
specialization. Homogeneous ternary-cubic factorial results are not the
same assertion as arbitrary cubic exponential nonvanishing on a triangle.

Papers II--III concern relations among compact affine-line periods, and
Paper IV concerns relative functional periods with its stated constant
field. They can organize face relations and relative systems, but do not
automatically give a numerical theorem for arbitrary simplex chains.
The candidate multicolumn workstream addresses another issue, namely
actual realization of affine orbit equations, and is likewise separate.

\paragraph{Reproducibility and status.}
The polynomial identity in \cref{prop:counterexample} and the simplex
Jacobians used in Paper VI are checked by
\texttt{scripts/check\_comparison\_factorial.py}. Such finite checks do not
prove the holonomic or transcendence inputs. This document preserves
proved reductions and precise open targets, not a completed proof of
\cref{conj:simplex}.

\begin{thebibliography}{9}
\raggedright
\bibitem{Beukers} F. Beukers, \emph{A refined version of the
Siegel--Shidlovskii theorem}, Ann. of Math. (2) \textbf{163} (2006), 369--379.
\url{https://annals.math.princeton.edu/wp-content/uploads/annals-v163-n1-p08.pdf}.
\bibitem{AR} B. Adamczewski and T. Rivoal, \emph{Exceptional values of
$E$-functions at algebraic points}, Bull. Lond. Math. Soc. \textbf{50}
(2018), 697--708. \url{https://adamczewski.perso.math.cnrs.fr/E-functions.pdf}.
\bibitem{Zeilberger} D. Zeilberger, \emph{A holonomic systems approach to
special functions identities}, J. Comput. Appl. Math. \textbf{32} (1990),
321--368.
\url{https://sites.math.rutgers.edu/~zeilberg/mamarimY/Zeilberger_y1991_p321.pdf}.
\bibitem{PaperI} C. D. Long and A.-A. Le Blanc, \emph{Algebraic Exponential
Integrals: Transcendence and Linear Relations}, Paper I, working manuscript.
\bibitem{PaperV} C. D. Long and A.-A. Le Blanc, \emph{Relative Exponential
Periods on Polynomial Chains: Algebraic Independence over Fixed Boundary
Fields}, Paper V, September 2026, revised October 2026.
\bibitem{PaperVI} C. D. Long and A.-A. Le Blanc, \emph{Exponential Periods
and Factorial Moments}, Paper VI, October 2026, working manuscript.
\end{thebibliography}
\end{document}
